\section{Inverting on one common domain}
\label{sec:stable-inverse}
\label{sec:grn-fixed-phase-green}

A tame estimate with derivative loss is not a bounded inverse on a
single Banach space. It cannot by itself justify a Neumann series.
We first identify a source space on which the \emph{relative}
operator is bounded. The extra endpoint factor in its coefficients
will compensate for the coefficient-weight loss of the reference
inverse. Only after that Banach-space inversion do we recover all
derivatives. This order of argument keeps one neighborhood for the
entire nonlinear problem.

\subsection{Sources, scalars and the domain}

Use the canonical cylinders of Section~\ref{sec:fixed-geometry} and
retain the physical sign of the equation: on a positive cylinder its
canonical operator is $-P_{\mathbf x}$. Let $h_{\mathbf x}$ denote the
height factor in physical coordinates, $\Psi_{\mathbf x}$ the chart,
and $J_{\mathbf x}$ its Jacobian. Define
\[
 \mathscr C_{\mathbf x}u=(u/h_{\mathbf x})\circ\Psi_{\mathbf x}^{-1},
 \qquad
 \mathscr S_{\mathbf x}f=(h_{\mathbf x}f/|J_{\mathbf x}|)
                                      \circ\Psi_{\mathbf x}^{-1},
 \qquad
 \mathbb L_{\mathbf x}=
 \mathscr S_{\mathbf x}\widehat L_{\mathbf x}\mathscr C_{\mathbf x}^{-1}.
\]
These preserve the pairing with $dr\,d\theta$. For real $t$ and
angular order $s\ge0$ use the volume-source norms
\begin{equation}\label{eq:vol-source-spaces}
 \|F\|_{t,s}^2
 =\sum_m M_m^{2t}H_m^{-2\beta}
                  \|F_m\|_{L^2_{\mathrm{rad}}H^s_\varphi}^2,
 \qquad F\in\mathbb V_{t,s}.
\end{equation}
The index includes the two halves of every negative component and
the positive components. No matching or integral condition is imposed
on these sources. For inverse outputs, the corresponding volume norm is supplemented
by the absolute values of the three scalar unknowns introduced below;
these scalar coordinates are assigned to the finite middle.

The charts agree with the reference on $J_*$, so the second-mode
normalization has the fixed form $\ell_*^c W$.
Choose a compact radial $\psi$ in $J_*$ with $\mathsf Q\psi=1$ and
put $c_* =\mathscr S_*\psi=\mathscr S_{\mathbf x}\psi$.
A physical height is represented by
\[
 \eta=\mathscr C_{\mathbf x}^{-1}W+a_L\chi_L,\qquad
 \ell_*^cW=0.
\]
Its left endpoint constant is $a_L$; the right one is zero. Define
\[
 k_{\mathbf x}=D_\nu\mathcal F_q(w,u),\quad
 k_{\mathbf x}^c=\mathscr S_{\mathbf x}k_{\mathbf x},\quad
 s_{\mathbf x}=\mathscr S_{\mathbf x}\widehat L_{\mathbf x}\chi_L.
\]
Here $D_\nu$ differentiates the true residual with $u,q$ fixed,
including both its metric and reference height. The augmented operator is
\begin{equation}\label{eq:vol-abstract-complete-operator}
 \mathbb P_{\mathbf x}(W,a_L,\mu,\lambda)
 =\mathbb L_{\mathbf x}W+a_Ls_{\mathbf x}
                       +\mu k_{\mathbf x}^c+\lambda c_*.
\end{equation}
The parameter $\mu$ removes the second-mode obstruction. The additional
$\lambda$ removes the integral obstruction for arbitrary sources;
it will disappear for the original compatible problem.

At the reference let $\mathfrak g_* =\mathscr C_*1$ and
$\mathsf Q_*F=\langle\mathfrak g_*,F\rangle$.
The coefficient bounds make this continuous and
$\mathsf Q_*c_*=1$. Set $\lambda=\mathsf Q_*F$ and apply the
compatible reference inverse to $F-\lambda c_*$. This constructs a
two-sided augmented inverse $\mathbb G_*$.
In the zero angular mode the construction can be seen directly:
with $f_* =\mathscr S_*^{-1}F$,
\[
 p_0(r)=\int_{r_-}^r(\av{f_*}-\lambda\psi)(t)\,dt,
 \qquad
 u_0(r)=-\int_r^{r_+}\frac{p_0(t)}{w_*(t)}\,dt.
\]
The first integrand has total integral zero. At either endpoint use
the primitive starting there; division by $w_*$ costs a coefficient
factor, not an endpoint exponent. Subtract the left endpoint value
times $\chi_L$. This yields the required field, endpoint scalar and
normalization. It also proves uniqueness in the zero mode.

\subsection{Why two endpoint exponents are needed}

The reference construction gives one natural endpoint inverse for
every exponent in a compact subinterval of $(0,1)$
(Lemma~\ref{lem:radial-uniform-exponents}). Choose
$0<\epsilon<\min(\beta,1-\beta)$.
The inverses at $\beta-\epsilon$ and $\beta+\epsilon$ agree on compact
smooth sources: their difference satisfies the homogeneous equation,
the normalization and the natural endpoint conditions. The reference
uniqueness, together with the zero-mode formula above, makes it zero.
Hence both estimates apply to the same $\mathbb G_*$.

\begin{lemma}[A fixed coefficient loss]
\label{lem:two-weights-closed-domain}
For a fixed $L$, every real $t$ and every $s\ge0$,
\begin{equation}\label{eq:reference-negative-material-extension}
 \|\mathbb G_*F\|^{\mathrm{vol}}_{t-L,s}
                       \le C_{t,s}\|F\|_{t,s}.
\end{equation}
The first normal derivative has the same bound with at most one
additional angular source derivative. For a fixed $b\ge2$,
\begin{equation}\label{eq:fixed-base-graph-domain}
 \mathbb D_*:=\mathbb G_*(\mathbb V_{0,b}),\qquad
 \|v\|_{\mathbb D_*}=\|\mathbb P_*v\|_{0,b}
\end{equation}
is a complete closed graph domain. It includes every smooth admissible
field with the stated transmissions, endpoint conditions and slice,
whose reference source belongs to $\mathbb V_{0,b}$.
\end{lemma}

\begin{proof}
Write $\bar S_m=S_m-S_N$ on the tails and $\bar S_m=0$ on the
finite middle. Let $\pi_k,\iota_m$ be restriction to an output piece
and extension of a source on a piece by zero. The reference estimate
at exponent $\gamma$ gives
\[
 \|\pi_k\mathbb G_*\iota_m\|_{s\to s}
 \le C_s\mathcal E_k^{C_s}\mathcal E_m^{C_s}
                        e^{-\gamma(\bar S_k-\bar S_m)}.
\]
Conjugate by the $\beta$ weight. Choose $\gamma=\beta+\epsilon$
when $\bar S_k\ge\bar S_m$ and $\gamma=\beta-\epsilon$ otherwise.
The resulting block $K_{km}$ satisfies
\begin{equation}\label{eq:global-two-weight-block-bound}
 \|K_{km}\|_{s\to s}
 \le C_s\mathcal E_k^{C_s}\mathcal E_m^{C_s}
                    e^{-\epsilon|\bar S_k-\bar S_m|}.
\end{equation}
There is no angular-order loss in this volume estimate. The first
normal estimate of the radial construction adds at most one angular
derivative and otherwise has the same block argument.

For each $\delta>0$ the enlarged weight satisfies
\[
 |\log M_k-\log M_m|
 \le\delta|\bar S_k-\bar S_m|+C_\delta.
\]
This follows from the derivative of
$n^{5/2}\log^2(e+n)$ with respect to $2n^3/3$ tending to zero.
Using $\mathcal E^{C_s}\le C_sM$, we find, for $L\ge2$,
\[
 M_k^{t-L}\|K_{km}\|M_m^{-t}
 \le C_sM_k^{2-L}(M_k/M_m)^{t-1}
                         e^{-\epsilon|\bar S_k-\bar S_m|}
 \le C_{t,s}e^{-(\epsilon/2)|\bar S_k-\bar S_m|}.
\]
Choose $\delta$ so that $|t-1|\delta\le\epsilon/2$.
The last matrix has bounded row and column sums: the $S_n$ are
uniformly separated and there are only finitely many pieces at each
index. The $\ell^2$ Schur estimate proves
\eqref{eq:reference-negative-material-extension}.

The definition of $\mathbb D_*$ makes it complete. Approximate a
source by compact smooth sources. The estimate gives local volume
convergence and the equation passes to distributions, with weak
transmissions. If $v_j\to v$ in the ambient volume space and
$\mathbb P_*v_j\to F$ in the source space, then
$v_j=\mathbb G_*\mathbb P_*v_j\to\mathbb G_*F$, so
$v=\mathbb G_*F$; the graph is closed. First-normal estimates identify
the local weak traces. Finally, any smooth admissible $y$ with source
$F$ differs from $\mathbb G_*F$ by a homogeneous normalized natural
endpoint solution, hence is equal to it. This proves the stated
exhaustion of admissible smooth fields.
\end{proof}

\begin{figure}[htbp]
\centering
\begin{tikzpicture}[x=.92cm,y=.9cm,font=\small,>=Stealth]
\draw[->] (0,0)--(6.2,0) node[right] {source depth $\bar S_m$};
\draw[->] (0,0)--(0,4.9) node[above] {output depth $\bar S_k$};
\fill[blue!7] (0,0)--(4.5,4.5)--(0,4.5)--cycle;
\fill[orange!9] (0,0)--(5.7,0)--(5.7,4.5)--(4.5,4.5)--cycle;
\draw[thick] (0,0)--(4.5,4.5);
\node at (1.45,3.4) {$\gamma=\beta+\epsilon$};
\node at (4,1.2) {$\gamma=\beta-\epsilon$};
\node[align=left,anchor=west] at (6.55,2.8)
 {$e^{-\epsilon|\bar S_k-\bar S_m|}$\\[3pt]
  controls transfer between pieces.};
\end{tikzpicture}
\caption{The same reference inverse is estimated with two endpoint
weights. Each weight is used on one side of the diagonal. Their minimum
gives decay in the distance between source and output depths.}
\label{fig:two-weights}
\end{figure}

\subsection{The estimate that explains cubic matching}

On a normalized positive cylinder consider
\[
 P_ju=-(au')'+j^2cu+Vu=F,
 \qquad a\ge\mathcal E_m^{-C},\quad
 c\ge\mathcal E_m^{-C}d,\quad d=y(1-y).
\]
Coefficient derivatives through order two, and $a^{-1}$, are bounded
by $\mathcal E_m^{C_1}$. Boundary values are arbitrary. We claim
\begin{equation}\label{eq:vol-arbitrary-boundary-graph-bound}
 \begin{split}
 &\|du'\|_2+\|\av j d^{3/2}u\|_2+\|d^2u''\|_2
   +\|\av j^2d^3u\|_2+\|\av j d^{5/2}u'\|_2\\
 &\hspace{30mm}\le\mathcal E_m^C(\|F\|_2+\|u\|_2),
 \end{split}
\end{equation}
where $\av j=(1+j^2)^{1/2}$ in this display.
For fixed $j$, a distributional $L^2$ solution with $F\in L^2$
is $H^2$ up to the endpoints by the regular ordinary differential
equation. Uniform bounds come from the following calculation, not
from that preliminary regularity statement.

Multiply by $d^2\bar u$ and integrate. Since $d=0$ at the ends,
\[
 \int a|(du)'|^2+j^2\int cd^2|u|^2
 =\operatorname{Re}\langle F,d^2u\rangle
       +\int a|d'|^2|u|^2-\int Vd^2|u|^2.
\]
This bounds $du'$ and $jd^{3/2}u$. Now put $v=d^2u$; it has
zero trace and
\[
 P_jv=d^2F-2a(d^2)'u'-(a(d^2)')'u.
\]
The right side is controlled by the first estimate. With $q_0=c/a$,
the same is true of $F_1=-v''+j^2q_0v$. Squaring and integrating
the cross term twice gives
\[
 \|v''\|_2^2+j^4\|q_0v\|_2^2+2j^2\int q_0|v'|^2
 =\|F_1\|_2^2+j^2\int q_0''|v|^2.
\]
The final integral is controlled because $d^4\le Cd^3$.
Use $q_0\ge\mathcal E_m^{-C}d$ and expand $v',v''$ to obtain
\eqref{eq:vol-arbitrary-boundary-graph-bound}.

Every second-order term in the perturbation
\eqref{eq:readable-positive-difference} is now controlled:
$d^3u_{yy}$ by $d^2u''$, $d^3u_{y\varphi}$ by
$\av j d^{5/2}u'$, and $d^3u_{\varphi\varphi}$ by
$\av j^2d^3u$. A derivative on the coefficient leaves only terms
controlled by the first two norms. Fourier convolution in $\varphi$
therefore gives, for $s\ge b$,
\begin{equation}\label{eq:readable-positive-relative-bound}
 \|\delta P_{\mathbf x}u\|_s
 \le C_sH_m^\beta\mathcal E_m^{C_s}
 \left[\mathfrak n_{b_0,m}(\|F\|_s+\|u\|_s)
       +\mathfrak n_{s+d_*,m}(\|F\|_b+\|u\|_b)\right].
\end{equation}
This is the reason for matching the principal density to order two.
It controls a second-order perturbation without imposing artificial
boundary conditions on the reference solution.

We also use an elementary first-normal estimate on every normalized
piece. Its radial equation has the form
$Au''+Bu'+(j^2C+D_0)u=F$ with $A$ bounded away from zero and all
coefficients bounded by fixed powers of $\mathcal E_m$.
The equation and
$\|u'\|_2^2\le C\|u\|_2(\|u''\|_2+\|u\|_2)$ give
\begin{equation}\label{eq:vol-reference-first-normal-bound}
 \|u'\|_2\le\mathcal E_m^C(\|F\|_2+\av j\|u\|_2),
 \qquad
 \|u''\|_2\le\mathcal E_m^C(\|F\|_2+\av j^2\|u\|_2).
\end{equation}
Young's inequality absorbs the intermediate $u'$ term. Density
extends the estimates to volume solutions. Neither estimate assumes
zero boundary data.

\subsection{The bounded relative operator}

Define on compact smooth sources
\[
 \mathscr R_{\mathbf x}=(\mathbb P_{\mathbf x}-\mathbb P_*)\mathbb G_*.
\]
If $\mathbb G_*F=(W,a_L,\mu,\lambda)$, then
\begin{equation}\label{eq:vol-normalized-relative-identity}
 \mathscr R_{\mathbf x}F
 =\delta\mathbb L_{\mathbf x}W
      +a_L(s_{\mathbf x}-s_*)+\mu(k_{\mathbf x}^c-k_*^c).
\end{equation}
The last scalar has no difference. On a negative piece only the
canonical potential varies. On a positive piece use
\eqref{eq:readable-positive-relative-bound}, putting its compact scalar
sources into $F$. Every coefficient difference at the output piece
$k$ contains $H_k^\beta$.

Here is the actual base-space calculation. Combine that factor with
\eqref{eq:global-two-weight-block-bound}. For low coefficients the
norm of the block from $m$ to $k$ is bounded by
\begin{equation}\label{eq:rewrite-relative-block}
 C\mathfrak n_{b_0}\,M_k^CM_m^C
            e^{-\beta\bar S_k-\epsilon|\bar S_k-\bar S_m|}.
\end{equation}
For every fixed $C$ and $\delta>0$,
$M_j^C\le C_\delta e^{\delta\bar S_j}$, while
$\bar S_m\le\bar S_k+|\bar S_k-\bar S_m|$.
Choose $\delta<\min(\beta/4,\epsilon/2)$. The matrix in
\eqref{eq:rewrite-relative-block} is bounded by
\[
 C'\mathfrak n_{b_0}
       e^{-(\beta/2)\bar S_k-(\epsilon/2)|\bar S_k-\bar S_m|}.
\]
Its row and column sums are uniformly bounded. Thus
\begin{equation}\label{eq:rewrite-base-relative}
 \|\mathscr R_{\mathbf x}\|_{\mathbb V_{0,b}\to\mathbb V_{0,b}}
                         \le C\mathfrak n_{b_0}<\tfrac12
\end{equation}
after increasing $b_0$ to cover the finite low derivatives and
then decreasing $\epsilon_0$. This is a bounded operator on precisely
one source space; no identification of a derivative-losing inverse
with a bounded endomorphism has been made.

For higher orders keep one coefficient factor at high order in the
product estimate. The same summation proves, for $t\ge0$, $s\ge b$,
\begin{equation}\label{eq:compact-volume-relative-bound}
 \begin{split}
 \|\mathscr R_{\mathbf x}F\|_{t,s}
 +\|[\Lambda,\mathscr R_{\mathbf x}]F\|_{t,s}
 \le C_{t,s}\bigl[\epsilon_0\|F\|_{0,s}
             +\mathfrak n_{t+s+d_*}\|F\|_{0,b}\bigr],
 \end{split}
\end{equation}
where $\Lambda=|D_\varphi|$. Assign frequency zero to $a_L,\lambda$
and frequency two to $\mu$. For multiplication by a coefficient,
the commutator Fourier matrix gains $|k|-|j|$, bounded by $|k-j|$;
one additional coefficient derivative therefore proves the same bound.
The scalar commutators are retained:
$\Lambda(s_{\mathbf x}-s_*)a_L$ and
$(\Lambda-2)(k_{\mathbf x}^c-k_*^c)\mu$.

Approximation by compact smooth sources identifies this extension
with the differential difference in distributions on $\mathbb D_*$.
The canonical transmissions are fixed, hence
\begin{equation}\label{eq:current-fixed-domain-factorization}
 \mathbb P_{\mathbf x}=(I+\mathscr R_{\mathbf x})\mathbb P_*
                         \quad\text{on }\mathbb D_*.
\end{equation}
This is the desired comparison on a common domain.

\subsection{Higher derivatives from the one base inverse}

By \eqref{eq:rewrite-base-relative}, set
\[
 \widetilde F=(I+\mathscr R_{\mathbf x})^{-1}F,
 \qquad v=\mathbb G_*\widetilde F.
\]
The inverse in the first formula is taken only on $\mathbb V_{0,b}$.
We now explain why this same solution has all the required regularity.

First recover angular derivatives. If $C=[\Lambda,\mathscr R]$, then
for $U_\delta=e^{i\delta\Lambda}$,
\[
 [U_\delta,\mathscr R]
 =i\int_0^\delta U_{\delta-t}CU_t\,dt.
\]
Apply this to $(I+\mathscr R)\widetilde F=F$. The base inverse
bounds the difference quotients of $\widetilde F$ whenever those
of $F$ are bounded. Fatou's lemma on Fourier coefficients gives the
first derivative; iteration gives all integer orders. Quantitatively,
the commuted equation at order $N$ contains
$\sum_{k=0}^{N-1}\Lambda^kC\Lambda^{N-1-k}\widetilde F$.
Interpolate the bound for $C$ at the base and at order $N-1$, treating
the high coefficient size as a weight in the source Hilbert norm.
Each term is then bounded by
\[
 C_N\bigl(\|\widetilde F\|_{N-1}
 +(1+\mathfrak n_{N+d_*})\|\widetilde F\|_b\bigr).
\]
Interpolation and Young's inequality absorb the first term into
half the order-$N$ norm plus a base norm. This proves the one-high-factor
estimate without smallness of high coefficient norms. Real orders
follow by Hilbert interpolation, or may be replaced by their ceilings.

Next use the already solved identity
$\widetilde F=F-\mathscr R_{\mathbf x}\widetilde F$.
Estimate its last term by \eqref{eq:compact-volume-relative-bound}.
It belongs to every coefficient-weight order once angular regularity
is known. Thus
\begin{equation}\label{eq:full-source-tame-resolvent}
 \|\widetilde F\|_{t,s}
 \le C_{t,s}\bigl[\|F\|_{t,s}
       +(1+\mathfrak n_{t+s+d_*})\|F\|_{0,b}\bigr].
\end{equation}
No second inversion is required.

\paragraph{The normal recursion uses the original source.}
Estimate \eqref{eq:full-source-tame-resolvent} controls angular
derivatives of $\widetilde F$. We do not assume normal derivatives
of that corrected source. Use instead the current equation
$\mathbb P_{\mathbf x}v=F$, and subtract its three scalar columns:
\[
 F^\circ=F-a_Ls_{\mathbf x}-\mu k_{\mathbf x}^c-\lambda c_*.
\]
Rescale each negative half to a unit radial interval, denoted by $y$;
the positive intervals already have this normalization. All rescaling
factors and their reciprocals are reference coefficient factors.
Solving the equation for its nonvanishing normal coefficient gives
\begin{equation}\label{eq:rewrite-normal-recursion-base}
 \partial_y^2W=\sum_{(i,j)\in I}A_{ij}\partial_y^i
                      \partial_\varphi^jW+B F^\circ,\qquad
 I=\{(1,1),(0,2),(1,0),(0,1),(0,0)\}.
\end{equation}
On a negative half the unscaled equation is
$W_{\xi\xi}=\xi W_{\varphi\varphi}-VW+F^\circ$.
On a positive piece the normal coefficient has a lower bound
$\mathcal E_m^{-C}$ on the common neighborhood. Thus all the
$A_{ij},B$ obey the coefficient calculus of
Proposition~\ref{prop:uniform-physical-geometry}, with a fixed
state-derivative loss $d_{\rm geo}$.

For $a\ge2$ the actual mixed recursion is
\begin{equation}\label{eq:rewrite-mixed-normal-recursion}
 \begin{split}
 \partial_y^a\partial_\varphi^b W
 ={}&\sum_{(i,j)\in I}\sum_{p=0}^{a-2}\sum_{q=0}^{b}
       \binom{a-2}{p}\binom bq
       (\partial_y^p\partial_\varphi^q A_{ij})
       \partial_y^{a-2-p+i}\partial_\varphi^{b-q+j}W\\
 &+\partial_y^{a-2}\partial_\varphi^b(BF^\circ).
 \end{split}
\end{equation}
Every field on the right has at most $a-1$ normal derivatives
and total order at most $a+b$. A derivative falling on a
coefficient lowers the remaining field order by that amount.
Consequently the expansion terminates in angular derivatives of
$W$, angular derivatives of $W_y$, and mixed derivatives of
$F^\circ$. At total order $s$ their respective orders are at most
$s$, $s-1$, and $s-2$.
Starting with the already established angular and first-normal
regularity, this identity in distributions proves the existence of
the mixed derivatives successively.

Here is the quantitative form of this recursion. Fix a base order
$b_1\ge4$, and let $\mathfrak n_{r,m}$ be the localized displacement
seminorm on the bounded-overlap neighborhoods used in
Section~\ref{sec:fixed-geometry}. For $s\ge b_1+2$ put
\[
 \mathcal A_{s,m}
 =\|W_m\|_{L_y^2H_\varphi^s}
   +\|\partial_yW_m\|_{L_y^2H_\varphi^{s-1}}
   +\|F_m^\circ\|_{H^{s-2}_{y,\varphi}}.
\]
The finite expansion of \eqref{eq:rewrite-mixed-normal-recursion}
gives
\begin{equation}\label{eq:rewrite-local-mixed-bound}
 \|W_m\|_{H^s_{y,\varphi}}
 \le C_s\mathcal E_m^{C_s}
       \left[\mathcal A_{s,m}
       +(1+\mathfrak n_{s+d_{\rm geo},m})\mathcal A_{b_1,m}\right].
\end{equation}
To justify the tame form, keep the derivative count in each
expanded term: the sum of coefficient derivative orders and the
remaining field or source order is at most $s$. Interpolate each
intermediate coefficient norm between its fixed low norm and its
order-$s+d_{\rm geo}$ norm, and each remaining field norm between
the corresponding entries of $\mathcal A_{b_1,m}$ and
$\mathcal A_{s,m}$. The interpolation exponents have sum at most
one. The weighted arithmetic--geometric mean inequality therefore
places at most one factor at high order. Sobolev embedding uses only
the fixed base order $b_1$. Reciprocal normal coefficients have the
same property by their positive lower bound. This proves
\eqref{eq:rewrite-local-mixed-bound}; it does not require smallness
of any order-$s$ state norm.

The reference factors in that finite expansion may have an exponent
$C_s$ growing with $s$. Absorb their \emph{complete product} only
after the expansion: $\mathcal E_m^{C_s}\le C_sM_m$ costs one
coefficient-weight order. Paying for each recursion separately
would obscure this fixed loss.
For precision, write
\[
 \bigl(\|Z\|_{t,s}^{\rm mix}\bigr)^2
 =\sum_m M_m^{2t}H_m^{-2\beta}
       \sum_{a+b\le s}\|\partial_y^a\partial_\varphi^b Z_m\|_2^2.
\]
The zeroth-normal entry of $\mathcal A_{s,m}$ is controlled by
\eqref{eq:reference-negative-material-extension}. For its
first-normal entry, use the first-normal reference bound at angular
order $s-1$: its one extra angular derivative asks for precisely
order $s$ of $\widetilde F$. Summing
\eqref{eq:rewrite-local-mixed-bound} therefore yields
\begin{equation}\label{eq:rewrite-global-mixed-bound}
 \begin{split}
 \|W\|_{t,s}^{\rm mix}
 \le C_{t,s}\bigl[&
   \|\widetilde F\|_{t+L+1,s}
   +\|F\|_{t+1,s-2}^{\rm mix}\\
 &+(1+\mathfrak n_{t+s+D})\|F\|_{0,b_1}^{\rm mix}\bigr].
 \end{split}
\end{equation}
Here $D$ is fixed: it covers $L+1$, $d_{\rm geo}$ and the fixed
base and embedding orders. The scalar columns are compactly
supported and satisfy the same tame estimate; their scalars are
bounded by the base inverse. In the high-coefficient term the
$M^{t+L+1}$ weight is placed on the localized state norm, leaving
the low field in the reference output space of weight $-L$.
Choose the original $b_0$ large enough to control the finitely many
base norms used here. The finitely many orders below $b_1+2$ follow
from this same base calculation.

Apply \eqref{eq:full-source-tame-resolvent} to the first term in
\eqref{eq:rewrite-global-mixed-bound}. The weight and derivative
indices there have sum $t+s+L+1$; the source term has sum $t+s-1$.
Thus taking $t+s\le r$ costs a fixed number of grading orders,
independent of $r$. The physical coordinate and source transformations
have their already proved fixed tame losses. Increasing $d_*$
once to include those losses gives \eqref{eq:readable-physical-inverse-tame}.
Physical seam jets match by the argument at the end of
Section~\ref{sec:fixed-geometry}.

Undoing the transformations gives
\[
 \widehat L_{\mathbf x}\eta+k_{\mathbf x}\mu+\lambda\psi=f,
 \qquad \ell_*\eta=0.
\]
Endpoint fluxes vanish. Moreover $\mathsf Qk_{\mathbf x}=0$ by
differentiating the integral compatibility of the original residual.
Integration therefore gives $\lambda=\mathsf Qf$. For
$f\in\mathcal Y_\beta$ it is zero.

\begin{proposition}[The perturbed scalar inverse]
\label{prop:rewrite-volume-inverse}
On the single neighborhood fixed above there is a unique normalized
pair $G_{\mathbf x}f=(\eta,\mu)$ solving
$\widehat L_{\mathbf x}\eta+k_{\mathbf x}\mu=f$ for every
$f\in\mathcal Y_\beta$. The height belongs to
$\mathcal X_\beta\cap\ker\ell_*$ and has right endpoint value zero.
Its left endpoint value is free. For a fixed $d_*$,
\begin{equation}\label{eq:readable-physical-inverse-tame}
 \|G_{\mathbf x}f\|_s
 \le C_s\bigl[\|f\|_{s+d_*}
       +(1+\mathfrak n_{s+d_*})\|f\|_{b+d_*}\bigr].
\end{equation}
The inverse is smooth tame jointly in state and source.
\end{proposition}

\begin{proof}
Existence and the estimate were proved above. Factorization on
$\mathbb D_*$ gives both inverse identities, and its exhaustion
property gives uniqueness among smooth admissible fields. The
resolvent identity first proves continuity in the base space and
then, using the higher estimates, at every stronger order. The
coefficient family is smooth tame. Hamilton's inverse-family theorem
\cite[II, Theorem~3.1.1, p.~150]{Hamilton1982} now applies: a smooth
tame linear family with unique solutions and a continuous tame inverse
has a smooth tame inverse. Its uniqueness and continuity hypotheses
have thus been checked before invoking it.
\end{proof}
